\section{Partial kernels under a modification}
\label{sec:birational-partial-kernels}

A modification identifies holomorphic sections, but it changes the integrability condition that defines multiplier ideals; the correction is a twist by the relative canonical bundle, through the Jacobian. If the relative canonical bundle carries the metric determined by two volume forms, the same Jacobian factor makes the identification of section spaces an isometry. The partial kernel measures downstairs are then exactly the pushforwards of those upstairs, and the convergence theorem on compact K\"ahler manifolds descends to every compact complex manifold that admits a K\"ahler modification.

Let $\pi\colon Y\to X$ be a modification of connected compact complex manifolds. We use the exceptional set $S\subset X$, its inverse image $D=\pi^{-1}(S)\subset Y$, the relative canonical bundle $R_\pi$ and its Jacobian section $J_\pi$ from \cref{sec:psh-birational-ideals}, together with the transformation of multiplier ideals in \cref{lem:adjoint-ideal-transform} and the invariance of the multiplier-ideal order in \cref{lem:birational-ideal-order}.

\begin{proposition}\label{prop:birational-partial-envelopes}
Let $\theta$ be a smooth real $(1,1)$-form on $X$ and let $\varphi\in\PSH(X,\theta)$. For $v\in C(X)$ let $E_X(v)=E_{\theta,X}[\varphi]_{\mathcal I}(v)$ be the partial envelope with obstacle $v$ on the whole of $X$. Then
\[
 E_{\pi^*\theta,Y}[\varphi\circ\pi]_{\mathcal I}(v\circ\pi)
       =E_X(v)\circ\pi.
\]
Moreover, for every $u\in\PSH(X,\theta)$ the extended non-pluripolar measures satisfy
\[
 \pi_*\mathrm{MA}^{\mathrm{ext}}_{\pi^*\theta}(u\circ\pi)
       =\mathrm{MA}^{\mathrm{ext}}_\theta(u).
\]
\end{proposition}
\begin{proof}
By \cref{prop:psh-descent-modification}, $a\mapsto a\circ\pi$ is a bijection from $\PSH(X,\theta)$ onto $\PSH(Y,\pi^*\theta)$. For a continuous obstacle on the whole manifold, the quasi-everywhere and the pointwise obstacle conditions are equivalent, as recalled in \cref{sec:obstacle-envelopes}, and since $\pi$ is surjective, $a\leq v$ on $X$ if and only if $a\circ\pi\leq v\circ\pi$ on $Y$. Together with \cref{lem:birational-ideal-order}, this identifies the two families of competitors. By \cref{thm:hermitian-partial-i-envelope}, each partial envelope is the largest member of its family, which proves the envelope identity; no interchange of pullback and regularized supremum is needed.

The map $\pi$ is a biholomorphism from $Y\setminus D$ onto $X\setminus S$, so locality of the non-pluripolar product identifies the two measures there. Neither measure charges the analytic sets $D$ and $S$. Testing against nonnegative Borel functions, or exhausting the complements by increasing compact sets, proves the pushforward identity, including the case of infinite total mass.
\end{proof}

Choose smooth positive probability volumes $\nu_X$ on $X$ and $\nu_Y$ on $Y$. They determine a smooth Hermitian metric $h_R$ on $R_\pi$ such that
\begin{equation}\label{eq:jacobian-volume-metric}
 |J_\pi|_{h_R}^2\,\nu_Y=\pi^*\nu_X.
\end{equation}
Explicitly, let $f_X$ and $f_Y$ be the densities of $\nu_X$ and $\nu_Y$ in local coordinates. The local frame of $R_\pi$ given by the quotient of the coordinate canonical frames is given the squared norm $(f_X\circ\pi)/f_Y$. This function is smooth and strictly positive, also where $J_\pi$ vanishes, and the transformation rules of the coordinate densities show that it defines a metric. Equation~\eqref{eq:jacobian-volume-metric} is an identity of smooth top-degree forms; it does not define the pullback of an arbitrary measure.

\begin{proposition}\label{prop:birational-partial-kernel-measures}\index{Bergman kernel!modification}
Let the line bundle $L$ with metric $h$, the selecting potential $\varphi\in\QPSH(X)$ and the continuous weight $v\in C(X)$ be as in \cref{sec:partial-selected-space}, and use the volume $\nu_X$ on $X$. On $Y$, use the line bundles $\pi^*L^m\otimes R_\pi$ with the metrics $\pi^*h^m\otimes h_R$, the selecting potential $\varphi\circ\pi$, the continuous weight $v\circ\pi$ and the volume $\nu_Y$, and define the kernels and kernel measures as in \cref{sec:partial-selected-space}, with the fixed factor $(R_\pi,h_R)$ added. Then the identification $s\mapsto J_\pi\pi^*s$,
\[
 H^0(X,L^m\otimes\mathcal I(m\varphi))
 \simeq H^0(Y,\pi^*L^m\otimes R_\pi
                         \otimes\mathcal I(m\varphi\circ\pi)),
\]
is an isometry for the continuous weighted norms. The kernels and kernel measures on $Y$ and $X$ satisfy
\[
 B_{Y,m}=|J_\pi|_{h_R}^2(B_{X,m}\circ\pi),
 \qquad \pi_*\beta_{Y,m}^{(r)}=\beta_{X,m}^{(r)}
\]
for the same normalization integer $r$, also when the section spaces are zero.
\end{proposition}
\begin{proof}
The identification is \cref{lem:adjoint-ideal-transform} with $M=L^m$ and $a=m\varphi$. By \eqref{eq:jacobian-volume-metric} and the change of variables formula,
\[
 \int_Y|J_\pi\pi^*s|_{\pi^*h^m\otimes h_R}^2\mathrm{e}^{-mv\circ\pi}\,\mathrm{d}\nu_Y
 =\int_X|s|_{h^m}^2\mathrm{e}^{-mv}\,\mathrm{d}\nu_X.
\]
The selecting potential has already determined the section spaces and does not enter the norms. Hence the identification maps an orthonormal basis to an orthonormal basis, and summing the squared pointwise norms of the basis vectors gives the kernel formula, also at the zeros of $J_\pi$. Multiplying by $r!/m^r$ and integrating a continuous function pulled back from $X$ proves the measure formula.
\end{proof}

\Needspace{25\baselineskip}
\subsection{Convergence with a K\"ahler modification}

Through the Jacobian isometry, the equilibrium theorem on a compact K\"ahler manifold can be pushed down along a modification. The limit does not depend on the modification, because both the partial envelope and its non-pluripolar Monge--Amp\`ere measure descend.

\begin{theorem}\label{thm:fujiki-partial-kernel-limit}\index{Equilibrium measure!Kahler modification@K\"ahler modification}
Suppose that $X$ admits a modification from a compact K\"ahler manifold. Let $L\to X$ be a holomorphic line bundle with a smooth Hermitian metric $h$, put $\theta=c_1(L,h)$, and let $\varphi\in\PSH(X,\theta)$ satisfy
\[
 \mathrm{MA}_\theta(\varphi)(X)>0.
\]
Then, for every $v\in C(X)$ and every smooth positive probability volume $\nu_X$, the partial kernel measures on $X$ satisfy
\[
 \frac{n!}{m^n}B_{m,v,\varphi,\nu_X}\,\nu_X
       \rightharpoonup\mathrm{MA}_\theta(E_X(v)).
\]
The limit measure is finite and nonzero. Its total mass is
\[
 M=\lim_{m\to\infty}\frac{n!}{m^n}
                    h^0(X,L^m\otimes\mathcal I(m\varphi)).
\]
\end{theorem}
\begin{proof}
Choose a modification $\pi\colon Y\to X$ from a compact K\"ahler manifold $Y$, and a smooth positive probability volume $\nu_Y$ on $Y$. The pullback $\varphi\circ\pi$ belongs to $\PSH(Y,\pi^*\theta)$, and by \cref{prop:birational-partial-envelopes} its non-pluripolar mass equals the given positive mass of $\varphi$. This mass is finite, by the theory of closed classes on compact K\"ahler manifolds.

Apply \cite[Theorem~14.3.1]{SGPT} on $Y$, with the line bundle $\pi^*L$, the auxiliary twisting line bundle $R_\pi$, the continuous obstacle $v\circ\pi$ on $Y$ and the probability volume $\nu_Y$. Its Bernstein--Markov hypothesis follows from the proof of \cref{prop:hermitian-full-bm}, since a fixed smooth auxiliary bundle only changes the local constants. The multiplier ideals and the continuous weighted norms in that theorem agree exactly with those of \cref{prop:birational-partial-kernel-measures}. The theorem gives
\[
 \beta_{Y,m}^{(n)}\rightharpoonup
 \mathrm{MA}_{\pi^*\theta}
 \bigl(E_{\pi^*\theta,Y}[\varphi\circ\pi]_{\mathcal I}(v\circ\pi)\bigr),
\]
and the limit measure has finite positive mass.

Now push this convergence forward to $X$, testing against continuous functions on $X$. By the kernel identity of \cref{prop:birational-partial-kernel-measures} and the envelope and pushforward identities of \cref{prop:birational-partial-envelopes}, this gives the asserted limit on $X$. Testing the constant function one shows that $n!N_m/m^n$ converges to the total mass of the limit, which proves the formula for $M$.
\end{proof}

Since $M>0$, the sections of $L$ grow at the maximal rate $m^n$, so $X$ is Moishezon by \cref{cor:non-moishezon-section-vanishing}. A K\"ahler modification also forces the bounded mass property, by \cref{lem:kahler-modification-bounded-mass}; the Jacobian identities, by contrast, hold for every modification of compact complex manifolds. When the section spaces grow more slowly than $m^n$, the total mass identity shows that the kernel measures must instead be normalized by the growth degree of $N_m$, as in \cref{sec:partial-submersion}.

\Needspace{15\baselineskip}
\begin{corollary}\label{cor:fujiki-partial-determinant}
Under the hypotheses of \cref{thm:fujiki-partial-kernel-limit}, let $M$ be the limit of the normalized dimensions in that theorem. The determinant functionals $\mathcal F_m^{(n)}$ of \cref{sec:partial-determinant} converge pointwise on $C(X)$ to
\[
 \mathcal F(v)=\int_0^1\!\left(\int_X
       v\,\mathrm{d}\mathrm{MA}_\theta(E_X(tv))\right)\mathrm{d}t.
\]
This functional is concave and vanishes at zero, and for $v,w\in C(X,\mathbb R)$ and $c\in\mathbb R$,
\[
 \mathcal F(v+c)=\mathcal F(v)+Mc,\qquad
 |\mathcal F(v)-\mathcal F(w)|\leq M\|v-w\|_\infty.
\]
\end{corollary}
\begin{proof}
Integrate the first-variation formula of \cref{prop:partial-kernel-variation} along the segment $t\mapsto tv$, $0\leq t\leq1$. For each fixed $t$, the integrand converges by the theorem, and its absolute value is bounded by
\[
 \|v\|_\infty\frac{n!N_m}{m^n},
\]
which is bounded in $m$ because it converges to $M\|v\|_\infty$. Dominated convergence on $[0,1]$ gives the formula; the limit integrand is measurable in $t$ as the pointwise limit of continuous functions. Concavity, the translation identity and the Lipschitz bound pass from $\mathcal F_m^{(n)}$ to the pointwise limit, since $n!N_m/m^n\to M$.
\end{proof}

This argument uses only weak convergence of measures along a segment of weights; it gives no lower bound for the partial Fubini--Study potentials complementing \cref{thm:partial-fubini-study-upper}.
